\begin{table}[H]
\centering
\caption{Effect of Turnout Decline on Crime}
\begin{adjustbox}{max width=\textwidth}
\begin{threeparttable}
\begin{tabular}{lccccccc}
\toprule
& (1) & (2) & (3) & (4) & (5) & (6) & (7) \\
& Total & Police- & Non-police- & Drug & Burglary & Homicide & Domestic \\
& crime & detected & dependent & offenses & & & violence \\
\midrule
Turnout decline $\times$ Post & $-0.0037$ & $-0.0090$** & $-0.0042$ & $-0.0471$*** & $-0.0162$*** & $+0.0132$*** & $+0.0028$ \\
                         & ($0.0044$) & ($0.0040$) & ($0.0037$) & ($0.0133$) & ($0.0043$) & ($0.0043$) & ($0.0046$) \\
\midrule
Comuna FE & Yes & Yes & Yes & Yes & Yes & Yes & Yes \\
Year FE & Yes & Yes & Yes & Yes & Yes & Yes & Yes \\
Observations & 3,061 & 3,061 & 3,061 & 3,061 & 3,061 & 3,061 & 3,061 \\
$R^2$ & 0.987 & 0.986 & 0.985 & 0.865 & 0.969 & 0.762 & 0.984 \\
\bottomrule
\end{tabular}
\begin{tablenotes}[flushleft]
\footnotesize
\item Notes: * p$<$0.10, ** p$<$0.05, *** p$<$0.01. Standard errors clustered at the comuna level in parentheses. The dependent variable is the natural log of annual crime counts (plus one). Turnout decline is the percentage-point drop in turnout from the 2008 (compulsory) to the 2012 (voluntary) municipal election. Post equals one for years 2018--2024. Column (7) serves as a placebo: domestic violence is victim-reported and should not respond to changes in police patrol intensity. N = 3,061 across 343 comunas. High $R^2$ in columns (1)--(3) and (5) reflects 343 comuna FEs absorbing cross-sectional variation; drug offenses and homicide have lower $R^2$ due to greater idiosyncratic variation.
\end{tablenotes}
\end{threeparttable}
\end{adjustbox}
\label{tab:main}
\end{table}
